% ===================================================================
%  บทความ: นิวรอนเทียมและการเรียนรู้ของเครื่อง
%  วิชา GE931-1 · หัวข้อ ปัญญาประดิษฐ์
%  วิธีคอมไพล์บน Overleaf: Menu -> Compiler -> XeLaTeX
%  ถ้าฟอนต์ไทยไม่ขึ้น ให้เปลี่ยน \setmainfont เป็นฟอนต์ไทยที่มีในระบบ
%  เช่น "Noto Sans Thai", "Sarabun", "Laksaman", "TH Sarabun New"
% ===================================================================
\documentclass[12pt]{article}

\usepackage{fontspec}
\usepackage{polyglossia}
\setmainlanguage{thai}
\setotherlanguage{english}
% --- ฟอนต์ไทย (แก้ตรงนี้ได้ถ้าฟอนต์ไม่ขึ้น) ---
\setmainfont{Laksaman}[Script=Thai]

\usepackage{amsmath}
\usepackage{amssymb}
\usepackage{geometry}
\geometry{a4paper, margin=2.5cm}
\usepackage{graphicx}
\usepackage{hyperref}
\hypersetup{colorlinks=true, linkcolor=black, urlcolor=blue}

\title{\textbf{นิวรอนเทียมและการเรียนรู้ของเครื่อง}}
\author{G \\ \small วิชา GE931-1 \quad หัวข้อ: ปัญญาประดิษฐ์}
\date{ปีการศึกษา 2569 ภาคเรียนที่ 1}

\begin{document}
\maketitle

\begin{abstract}
บทความนี้อธิบายแนวคิดพื้นฐานของ \textit{นิวรอนเทียม} (artificial neuron)
ซึ่งเป็นหน่วยย่อยของโครงข่ายประสาทเทียม และกระบวนการ ``เรียนรู้'' ของเครื่อง
ผ่านการปรับค่าน้ำหนักเพื่อลดค่าความสูญเสีย (loss) โดยยกตัวอย่างฟังก์ชันกระตุ้น
แบบซิกมอยด์และฟังก์ชันสูญเสียแบบครอสเอนโทรปี เพื่อให้ผู้อ่านเห็นภาพว่า
ปัญญาประดิษฐ์สมัยใหม่ ``เรียนจากข้อมูล'' ได้อย่างไร
\end{abstract}

\section{บทนำ}
ปัญญาประดิษฐ์ (Artificial Intelligence) ในปัจจุบันส่วนใหญ่ขับเคลื่อนด้วย
การเรียนรู้ของเครื่อง (Machine Learning) แทนที่จะเขียนกฎเกณฑ์ทุกข้อด้วยมือ
เราป้อน \textbf{ข้อมูลตัวอย่าง} จำนวนมากให้ระบบ แล้วให้ระบบปรับตัวเองจนทำนายได้แม่นยำ
หน่วยพื้นฐานที่ทำให้สิ่งนี้เกิดขึ้นได้คือ ``นิวรอนเทียม''

\section{โมเดลของนิวรอน}
นิวรอนเทียมรับค่าอินพุต $\mathbf{x}=(x_1,\dots,x_d)$ คูณด้วยน้ำหนัก
$\mathbf{w}$ บวกค่าไบแอส $b$ แล้วส่งผ่านฟังก์ชันกระตุ้น $\sigma$
ได้ผลลัพธ์ตามสมการที่~\eqref{eq:neuron}

\begin{equation}
\hat{y} \;=\; \sigma\!\left(\mathbf{w}^{\top}\mathbf{x}+b\right),
\qquad
\sigma(z)=\frac{1}{1+e^{-z}}
\label{eq:neuron}
\end{equation}

ฟังก์ชันซิกมอยด์ $\sigma(z)$ จะบีบค่าให้อยู่ในช่วง $(0,1)$
จึงตีความเป็น ``ความน่าจะเป็น'' ได้

\section{การเรียนรู้}
การเรียนรู้คือการหาน้ำหนัก $\mathbf{w},b$ ที่ทำให้ค่าความสูญเสียต่ำที่สุด
สำหรับปัญหาจำแนกสองกลุ่ม นิยมใช้ฟังก์ชันสูญเสียแบบครอสเอนโทรปี
ตามสมการที่~\eqref{eq:loss}

\begin{equation}
\mathcal{L} \;=\; -\frac{1}{N}\sum_{i=1}^{N}
\Big[\,y_i\log \hat{y}_i + (1-y_i)\log(1-\hat{y}_i)\,\Big]
\label{eq:loss}
\end{equation}

จากนั้นใช้วิธี \textit{gradient descent} ปรับน้ำหนักไปในทิศที่ทำให้
$\mathcal{L}$ ลดลงทีละน้อย จนกระทั่งโมเดลทำนายได้แม่นยำเพียงพอ

\section{สรุป}
นิวรอนเทียมเพียงหน่วยเดียวคือสมการเชิงเส้นที่ตามด้วยฟังก์ชันไม่เชิงเส้น
แต่เมื่อนำมาต่อกันเป็นชั้น ๆ จำนวนมากก็กลายเป็นโครงข่ายประสาทเชิงลึก
ที่สามารถเรียนรู้รูปแบบซับซ้อนจากข้อมูลได้ ซึ่งเป็นรากฐานของปัญญาประดิษฐ์ยุคใหม่

\begin{thebibliography}{9}
\bibitem{goodfellow}
Ian Goodfellow, Yoshua Bengio, and Aaron Courville.
\textit{Deep Learning}. MIT Press, 2016.

\bibitem{nielsen}
Michael Nielsen.
\textit{Neural Networks and Deep Learning}.
Determination Press, 2015.
\end{thebibliography}

\end{document}
